% Глава: 05-permissions
% Статус: наполнено (первый прогон).
\chapter{Права доступа и пользователи}
\label{ch:05-permissions}

\section{Пользователи и группы}

Linux --- многопользовательская система: каждый процесс работает от
имени \textbf{пользователя} и принадлежит \textbf{группе}. Каждому
пользователю присвоено числовое идентификатор (UID), каждой группе ---
GID. Текстовые имена хранятся в \fpath{/etc/passwd} (пользователи)
и \fpath{/etc/group} (группы).

На сервере \fpath{pglab.nsu.ru} вы --- обычный пользователь своей ВМ;
привилегии суперпользователя получают только администраторы.

Свой аккаунт и группы:
\begin{minted}[bgcolor=gray!6]{bash}
id                     # uid, gid и все группы текущего пользователя
whoami                 # имя текущего пользователя
groups                 # группы
getent passwd          # список пользователей, включая LDAP
\end{minted}

Пользователь с UID~0 --- \cmd{root} (администратор). Обычные
пользователи не могут менять системные файлы, пока не получат
временные привилегии.

\section{Права на файл}

У каждого файла (и каталога) есть \textbf{владелец},
\textbf{группа} и \textbf{права}. Права задаются для трёх субъектов:
\textbf{владелец} (owner), \textbf{группа} (group),
\textbf{остальные} (others). Для каждого субъекта возможны
три операции: \textbf{чтение} (r), \textbf{запись} (w),
\textbf{исполнение} (x).

Разбор записи \cmd{ls -l}:
\begin{minted}[bgcolor=gray!6]{bash}
-rw-r--r-- 1 alice alice 1234 Jan 1 12:00 notes.txt
drwxr-xr-x 2 alice devgrp    4096 Jan 1 12:00 project
\end{minted}

Первый символ --- тип файла: \cmd{-} обычный, \cmd{d} каталог,
\cmd{l} символическая ссылка. Затем три группы по три буквы:
владелец, группа, остальные. В примере владелец читает и пишет
\emph{и исполняет}, группа и остальные --- только читают.

Значение \textbf{исполнения} зависит от типа:
для файла --- можно запустить как программу; для каталога ---
можно войти внутрь (\cmd{cd}) и видеть содержимое.
Без \cmd{x} на каталоге нельзя перейти в него, даже если есть \cmd{r}.

\section{Числовая запись и \cmd{chmod}}

Права удобно задавать числом: для каждого субъекта
\cmd{r=4}, \cmd{w=2}, \cmd{x=1}, сумма --- одна цифра.
Например \cmd{7} (rwX), \cmd{5} (r-x), \cmd{6} (rw-), \cmd{4} (r--).
\begin{minted}[bgcolor=gray!6]{bash}
chmod 755 script.sh     # rwxr-xr-x
chmod 644 notes.txt     # rw-r--r--
chmod +x script.sh      # дать право выполнения всем
chmod u+w file.txt      # дать владельцу право записи
chmod a-x file.txt      # убрать право выполнения у всех
\end{minted}

\textbf{umask} --- маска, из которой система вычитает права
при создании файлов: это объясняет, почему новый файл
получает \cmd{644}, а новый каталог --- \cmd{755}.
\cmd{umask} без аргументов показывает текущую маску.

\section{Смена владельца и группы}

Собственником файлов можно стать только от имени \cmd{root}
(обычно через \cmd{sudo}, см.\ ниже).
\begin{minted}[bgcolor=gray!6]{bash}
sudo chown alice:devgrp file.txt   # сменить владельца и группу
sudo chown alice file.txt          # только владелец
chgrp devgrp file.txt              # группа, если вы в этой группе
\end{minted}

\section{\cmd{sudo}: привилегии администратора}

\cmd{sudo} выполняет команду \textbf{от имени другого пользователя}
(по умолчанию \cmd{root}) после проверки. Это основной способ
делать системные действия (установка пакетов, смена владельца):
\begin{minted}[bgcolor=gray!6]{bash}
sudo apt update            # выполнить команду от root
sudo -u postgres psql      # выполнить команду от пользователя postgres
sudo su -                  # получить оболочку root
\end{minted}

После первой проверки пароля в течение нескольких минут
\cmd{sudo} не спрашивает его повторно. Пользователь
попал в список, кому разрешён \cmd{sudo}, --- в
\fpath{/etc/sudoers} (редактируется через \cmd{visudo}).

\begin{warnbox}
Внутри сессии под \cmd{root} (после \cmd{sudo su -})
легко ошибиться с путём при \cmd{rm}. Для одноразовых
системных команд лучше \cmd{sudo <команда>},
чем входить в root-оболочку.
\end{warnbox}

\section{Частые специальные права (кратко)}

Существуют расширенные биты:
\textbf{setuid} и \textbf{setgid} (файл выполняется
с правами владельца/группы файла) и \textbf{sticky bit}
на каталоге (удалять файлы в нём может только владелец ---
так устроено \fpath{/tmp}). Они видны как \cmd{s}, \cmd{S},
\cmd{t} в записи \cmd{ls -l}. Для курса достаточно знать их
наличие и где они встречаются (например, \cmd{su}, \cmd{passwd}
имеют setuid).

\section{Типичные ошибки}
\begin{itemize}
  \item \emph{Permission denied} при запуске скрипта --- нет \cmd{x};
    дайте \cmd{chmod +x}.
  \item Нельзя открыть файл в каталоге --- нет \cmd{r} на каталоге
    \emph{или} нет \cmd{x} на нём (нужны оба).
  \item \cmd{chmod} не сработал --- вы не владелец и не \cmd{root}.
  \item Скрипт с \cmd{\#!/bin/bash} не запускается --- нет \cmd{x}
    или неверная строка интерпретатора в первой строке.
  \item Не работает \cmd{sudo} --- пользователь не в группе,
    которой разрешён sudo (проверьте с преподавателем).
\end{itemize}

\section{Связь с остальным курсом}
Права на файлы --- основа безопасности. Монтирование в Docker
(\nameref{ch:21-docker-data}) работает по той же модели:
что контейнеру доступно, определяется правами и точками монтирования.
Пакеты ставятся с правами \cmd{root}
(\nameref{ch:08-packages}), а работа под другим пользователем
(\cmd{sudo -u}) часто нужна при запуске сервисов.

\section{Задачи для самопроверки}
\begin{enumerate}
  \item Расшифруйте запись \fpath{drwxr-x---}: тип файла, права
    владельца, группы, остальных.
  \item Числовой записью задайте права «владелец читает/пишет,
    остальные --- только чтение». Какую команду вы выполнили?
  \item Создайте каталог, в который вам нельзя войти, и объясните
    по каким правам это произошло.
  \item Объясните разницу между \cmd{chmod u+w} и \cmd{chmod 666}.
  \item Как выполнить команду один раз от имени другого
    пользователя, не заходя в его оболочку?
\end{enumerate}
